\chapter{Wide local excision}
\index{Wide local excision}

\section*{Definition and Overview}
Wide local excision (WLE) is a definitive surgical procedure designed to remove a primary tumor or high-risk pre-malignant lesion along with a specified margin of surrounding healthy tissue. The fundamental objective of WLE is to achieve complete local control of the malignancy, verified histologically by the presence of microscopically clear surgical margins (R0 resection). By removing a peripheral boundary of healthy stroma, the surgeon targets subclinical, microscopic extensions of tumor cells that may have invaded the adjacent tissue but remain invisible to visual inspection or palpation.

WLE is widely utilized across multiple oncological disciplines, most notably in cutaneous oncology (for melanoma and non-melanoma skin cancers), breast surgery (where it is the cornerstone of breast-conserving therapy and referred to as a lumpectomy or segmentectomy), and soft tissue oncology (for sarcomas). The width of the surgical margin is dictated by evidence-based guidelines tailored to the tumor type, anatomical location, and staging parameters. For cutaneous melanoma, margins are defined by the Breslow thickness of the primary lesion, ranging from 5 to 10 millimeters for melanoma in situ, up to 2 centimeters for lesions thicker than 2.0 millimeters. In breast-conserving surgery, the standard is the ``no ink on tumor'' rule for invasive carcinoma, and a 2 millimeter margin for ductal carcinoma in situ (DCIS).

\section*{Epidemiology}
Because wide local excision is a surgical technique, its epidemiology directly reflects the incidence and demographics of the primary malignancies it treats. Globally, the demand for this procedure is driven by the rising incidence of skin and breast cancers.

Cutaneous melanoma represents a major public health challenge in light-skinned populations, with the highest incidence rates recorded in many countries and New Zealand, where the age-standardized rate exceeds 50 cases per 100,000 individuals annually. In the United States, melanoma is the fifth most common cancer, with a lifetime risk of approximately 1 in 38 for Caucasians. Non-melanoma skin cancers, specifically basal cell carcinoma (BCC) and squamous cell carcinoma (SCC), are the most common cancers globally, with millions of cases diagnosed annually.

Breast cancer is the most frequently diagnosed malignancy in women worldwide, with over 2 million new cases annually. In developed nations, breast-conserving surgery via WLE represents the primary treatment choice for 60\% to 70\% of early-stage breast cancer patients. Soft tissue sarcomas represent less than 1\% of adult malignancies, yet they represent a critical indication for WLE, typically affecting adults with a median age of 65 years.

\section*{Aetiology and Pathophysiology}
The necessity for a wide local excision is rooted in the infiltrative biology of malignant cells, which migrate along tissue planes and lymphatic channels.

Key pathophysiological mechanisms include:
\begin{itemize}
    \item \textbf{Ultraviolet (UV) Radiation-Induced Mutagenesis}: The primary driver for cutaneous malignancies is exposure to UV radiation. UVB radiation directly causes DNA damage by forming cyclobutane pyrimidine dimers, while UVA radiation acts indirectly via oxidative stress. These lead to signature mutations in critical genes, such as the \textit{BRAF} V600E mutation in 50\% of melanomas, activating the MAPK pathway, and \textit{TP53} mutations in SCC, disabling cell cycle arrest.
    \item \textbf{Microscopic Local Invasion}: Malignant cells undergo an epithelial-to-mesenchymal transition (EMT), upregulating matrix metalloproteinases (MMPs) that degrade the basement membrane. Individual tumor cells migrate into the surrounding stroma, extending millimeters beyond the macroscopic border of the tumor.
    \item \textbf{Lymphatic and Vascular Spread}: As tumors enlarge, localized hypoxia triggers the secretion of vascular endothelial growth factors (VEGF-A and VEGF-C), promoting angiogenesis and lymphangiogenesis. Tumor cells can enter these peritumoral channels, leading to satellite or in-transit metastases. Excision margins must encompass these local lymphatic routes.
    \item \textbf{Genetic Susceptibility}: Inherited mutations, such as \textit{BRCA1} or \textit{BRCA2} in breast cancer, impair homologous recombination, causing genomic instability. Germline mutations in \textit{CDKN2A} in familial melanoma impair cell cycle regulation, leading to aggressive, invasive phenotypes that require wide margins.
    \item \textbf{Tumor Microenvironment Interaction}: Surrounding stroma is remodeled by cancer-associated fibroblasts (CAFs), which deposit collagen tracks that facilitate tumor cell migration, requiring a wider excision to clear the invasive front.
\end{itemize}

\section*{Clinical Features}
The clinical presentation of patients requiring a wide local excision depends on the site and type of the primary neoplasm. Because WLE is a definitive secondary treatment, patients have typically undergone an initial diagnostic biopsy.

Key clinical features include:
\begin{itemize}
    \item \textbf{Cutaneous Melanoma}: Evaluated using the classic ABCDE criteria:
    \begin{itemize}
        \item \textit{Asymmetry}: One half of the lesion does not match the other.
        \item \textit{Border}: Irregular, notched, or blurred edges.
        \item \textit{Color}: Variations in color, including shades of brown, black, red, white, or blue.
        \item \textit{Diameter}: Usually greater than 6 millimeters.
        \item \textit{Evolving}: Changes in size, shape, color, or symptoms such as bleeding or itching.
    \end{itemize}
    \item \textbf{Non-Melanoma Skin Cancers}: BCC presents as a pearly, translucent nodule with telangiectasia and central ulceration (rodent ulcer). SCC presents as a firm, erythematous, hyperkeratotic plaque or nodule, often on sun-exposed skin.
    \item \textbf{Breast Malignancies}: Present as a hard, painless, irregular, and fixed breast lump. Associated signs include skin dimpling (peau d'orange), nipple retraction, bloody nipple discharge, or axillary lymphadenopathy.
    \item \textbf{Soft Tissue Sarcomas}: Present as a deep-seated, firm, painless mass in the extremities or trunk that progressively increases in size, often exceeding 5 centimeters at presentation.
\end{itemize}

\section*{Differential Diagnosis}
A range of benign and malignant lesions can mimic the clinical appearance of cancers requiring WLE. Differentiating these is critical to prevent unnecessary surgical morbidity.

Key differentials include:
\begin{itemize}
    \item \textbf{Benign Melanocytic Naevi}: Common moles that are symmetrical, uniform in color, and have regular borders. They lack cytologic atypia and pagetoid spread on histology.
    \item \textbf{Seborrheic Keratosis}: Common benign epidermal lesions with a ``stuck-on'' appearance, waxy texture, and horn pseudocysts, showing no dermal invasion.
    \item \textbf{Dermatofibroma}: A benign dermal nodule that exhibits the ``dimple sign'' upon lateral compression, differentiating it from nodular melanoma.
    \item \textbf{Dysplastic Naevi}: Atypical moles that display some clinical features of melanoma but lack malignant histopathological criteria.
    \item \textbf{Fibroadenoma}: The most common benign breast mass in young women, presenting as smooth, firm, and highly mobile (``breast mouse'').
    \item \textbf{Fibrocystic Changes}: Benign, cyclic breast changes that present with diffuse nodularity and tenderness, differentiated by imaging.
    \item \textbf{Lipoma}: A soft, mobile, subcutaneous mass of mature adipocytes, distinguished from liposarcoma which is deeper, firmer, and fixed.
    \item \textbf{Actinic Keratosis}: A scaly, premalignant lesion on sun-exposed skin that represents a precursor to SCC but lacks the invasive dermal component.
\end{itemize}

\section*{When to Seek Medical Care}
Individuals should consult a doctor if they detect any warning signs of malignancy, such as a changing mole, a non-healing skin sore, a new breast lump, or a rapidly expanding soft tissue mass.

Immediate emergency care must be sought for life-threatening symptoms, particularly post-operatively. These include:
\begin{itemize}
    \item Active, uncontrollable bleeding from the surgical wound.
    \item Sudden shortness of breath, severe chest pain, or coughing up blood, which may indicate a pulmonary embolism.
    \item Systemic signs of sepsis, such as a fever above 38.5 degrees Celsius, shaking chills, rapid heart rate, confusion, or cold, clammy skin.
    \item Severe, worsening pain or spreading redness around the incision site.
\end{itemize}
In the event of an emergency, contact services immediately:
\begin{itemize}
    \item In many countries, dial local emergency services.
    \item In the United States, dial 911.
    \item In the United Kingdom, dial 999.
    \item In Europe, dial 112.
\end{itemize}

\section*{Diagnosis and Investigations}
Before a wide local excision can be performed, a definitive tissue diagnosis must be established through investigations.

Key diagnostic steps include:
\begin{itemize}
    \item \textbf{Diagnostic Biopsy}:
    \begin{itemize}
        \item \textit{Skin}: An excisional biopsy with 1 to 3 millimeter margins is preferred to preserve lymphatic channels for sentinel node mapping.
        \item \textit{Breast}: Core needle biopsy (CNB) under ultrasound or mammographic guidance is mandatory to distinguish invasive cancer from DCIS.
        \item \textit{Sarcoma}: Core needle biopsy is preferred; the biopsy tract must be planned so it can be completely excised during the WLE.
    \end{itemize}
    \item \textbf{Histopathological Parameters}: The pathology report determines WLE margins based on Breslow thickness, mitotic rate, ulceration, and margin status of the biopsy specimen.
    \item \textbf{Imaging Studies}:
    \begin{itemize}
        \item \textit{Breast}: Mammography and ultrasound are standard. Breast MRI is used to define the extent of disease.
        \item \textit{Sarcoma}: MRI of the affected area is the gold standard to evaluate relationship to nerves, vessels, and bone.
        \item \textit{Staging}: CT scans or PET/CT are used for high-risk tumors to rule out distant metastases.
    \end{itemize}
    \item \textbf{Sentinel Lymph Node Biopsy (SLNB)}: Indicated for invasive melanoma (Breslow thickness greater than 1.0 millimeter, or greater than 0.8 millimeter with ulceration) and clinically node-negative invasive breast cancer. A radiotracer and blue dye are injected peri-tumorally to identify the first draining lymph node, which is excised to check for micrometastases.
\end{itemize}

\section*{Management}
The management of patients requiring WLE involves surgical excision, wound reconstruction, and often adjuvant therapy.

Key management components include:
\begin{itemize}
    \item \textbf{Surgical Excision}: Under appropriate anesthesia, WLE is performed with margins tailored to the lesion:
    \begin{itemize}
        \item \textit{Melanoma}: 5 to 10 millimeter margins for in situ; 1 centimeter for thickness less than or equal to 1.0 millimeter; 1 to 2 centimeters for thickness 1.01 to 2.0 millimeters; 2 centimeters for thickness greater than 2.0 millimeters.
        \item \textit{Breast Cancer}: Lumpectomy with macroscopic clear margins (``no ink on tumor'' for invasive disease; 2 millimeters for DCIS).
        \item \textit{Sarcoma}: 1 to 2 centimeter margins in three dimensions, extending to the deep fascia.
    \end{itemize}
    \item \textbf{Wound Reconstruction}: Closure is achieved based on defect size:
    \begin{itemize}
        \item \textit{Primary Closure}: Direct suturing for small defects.
        \item \textit{Skin Grafts}: Split-thickness (SSG) or full-thickness (FTSG) grafts harvested from donor sites for larger defects.
        \item \textit{Flaps}: Local tissue flaps or oncoplastic breast reconstruction to preserve vascular supply and cosmetics.
    \end{itemize}
    \item \textbf{Adjuvant Treatments}:
    \begin{itemize}
        \item \textit{Radiotherapy}: Standard after breast lumpectomy and sarcoma resection to reduce local recurrence.
        \item \textit{Systemic Therapy}: Adjuvant immunotherapy (e.g., Pembrolizumab) or targeted therapy for stage III melanoma. Endocrine therapy, chemotherapy, or HER2-targeted agents for breast cancer.
    \end{itemize}
\end{itemize}

\section*{Complications}
Wide local excision carries risks of both short-term (acute) and long-term (chronic) complications.

Potential complications include:
\begin{itemize}
    \item \textbf{Hemorrhage and Hematoma}: Bleeding into the wound bed, which may require surgical evacuation if large.
    \item \textbf{Surgical Site Infection (SSI)}: Bacterial infection presenting with erythema, warmth, pain, and purulent discharge, treated with antibiotics.
    \item \textbf{Wound Dehiscence}: Separation of the wound edges, often due to tension, poor blood supply, or infection.
    \item \textbf{Seroma}: Accumulation of serous fluid in the surgical cavity, requiring needle aspiration if symptomatic.
    \item \textbf{Graft or Flap Necrosis}: Tissue death in the reconstructed area due to compromised blood supply.
    \item \textbf{Chronic Pain and Sensory Loss}: Numbness, paresthesia, or neuropathic pain due to transected cutaneous nerves.
    \item \textbf{Lymphedema}: Chronic swelling of the limb, a risk when WLE is combined with sentinel lymph node biopsy or dissection.
    \item \textbf{Cosmetic Deformity}: Scarring, skin contour defects, or breast asymmetry, which may require secondary reconstruction.
\end{itemize}

\section*{Prognosis and Outcomes}
The prognosis following WLE is generally excellent for early-stage disease when clear margins (R0 resection) are achieved.

Key prognostic factors include:
\begin{itemize}
    \item \textbf{Local Control}: WLE achieves high local control rates. Melanoma recurrence is less than 3\% for thin tumors. Breast lumpectomy with radiotherapy yields a 10-year local recurrence rate of 5\% to 10\%, equivalent to mastectomy.
    \item \textbf{Margin Clearance}: Close or positive margins double the risk of local recurrence, necessitating re-excision.
    \item \textbf{Tumor Stage}: Survival is primarily dictated by stage. The 5-year survival for localized melanoma is greater than 98\%, dropping to 60\% to 70\% if regional lymph nodes are involved.
    \item \textbf{Recovery Timeline}: Suture removal occurs within 10 to 14 days. Complete tissue remodeling and nerve regeneration continue for 12 to 18 months.
\end{itemize}

\section*{Prevention and Self-Care}
Post-operative care and long-term prevention are crucial for optimal healing and detecting recurrence.

Self-care and prevention steps include:
\begin{itemize}
    \item \textbf{Wound Care and Hygiene}: Keeping the dressing dry for 24 to 48 hours, then washing gently. Avoid soaking in water. Apply sterile petroleum jelly to promote moist healing and reduce scarring.
    \item \textbf{Activity Management}: Restricting heavy lifting or strenuous activity that puts tension on the wound to prevent dehiscence.
    \item \textbf{Sun Safety}: For skin cancer prevention, apply SPF 50+ sunscreen, wear protective clothing, hats, and sunglasses, and avoid peak UV hours. Avoid tanning beds.
    \item \textbf{Self-Monitoring}: Perform monthly skin self-exams (using the ABCDE criteria) and regular breast self-exams to detect new or changing lesions.
    \item \textbf{Clinical Follow-Up}: Attend regular check-ups (every 3 to 6 months for early melanoma survivors) and annual screening mammograms for breast cancer survivors.
\end{itemize}

\section*{Sources and Further Reading}
\begin{itemize}
    \item Cancer Council Australia. \textit{Clinical Practice Guidelines for the Diagnosis and Management of Melanoma}. Available at: \url{https://www.cancer.org.au/clinical-guidelines/skin-cancer/melanoma} (Accessed: 9 August 2026).
    \item American Cancer Society. \textit{Treating Melanoma Skin Cancer}. Available at: \url{https://www.cancer.org/cancer/types/melanoma-skin-cancer/treating.html} (Accessed: 9 August 2026).
    \item National Comprehensive Cancer Network (NCCN). \textit{NCCN Clinical Practice Guidelines in Oncology: Breast Cancer and Melanoma}. Available at: \url{https://www.nccn.org/guidelines} (Accessed: 9 August 2026).
    \item Mayo Clinic. \textit{Lumpectomy: Wide Local Excision for Breast Cancer}. Available at: \url{https://www.mayoclinic.org/tests-procedures/lumpectomy/about/pac-20394650} (Accessed: 9 August 2026).
    \item Melanoma Institute Australia. \textit{Understanding Melanoma: Surgical Excision}. Available at: \url{https://www.melanoma.org.au/understanding-melanoma/stages-and-treatment/surgery} (Accessed: 9 August 2026).
    \item British Association of Dermatologists. \textit{Patient Information Leaflet: Wide Local Excision}. Available at: \url{https://www.bad.org.uk/patient-information-leaflets} (Accessed: 9 August 2026).
    \item World Health Organization (WHO). \textit{International Agency for Research on Cancer (IARC): Skin Cancer Prevention and Global Statistics}. Available at: \url{https://www.iarc.who.int} (Accessed: 9 August 2026).
\end{itemize}